\begin{itemize}[nosep,leftmargin=*]
\item \prop{infl}: \emph{Structural reasoning and conditional inference.}\\
  \bprint{INFL(PrtDriver,Problem1|PTROFFLINE)∧}\\
  \bprint{P(Problem1=No\_Output|PrtDriver=No)≥0.5}\\
  The printer driver still influences the output even when the offline status is known, and a missing driver makes failure likely.
\item \prop{map\_conf}: \emph{Explanation confidence.}\\
  \bprint{P(MAP(NetPrint,NetOK|Problem1=No\_Output)|}\\
  \bprint{Problem1=No\_Output)>=0.5-0.01}\\
  It accounts for at least $0.49$ of the posterior mass.
\item \prop{map\_update}: \emph{What-if analysis.}\\
  \bprint{P(MAP(NetPrint,NetOK|Problem1=No\_Output)[AppOK=Incorrect\_}\\
  \bprint{Corrupt↦0.5]↔MAP(NetPrint,NetOK|Problem1=No\_Output))=1}\\
  Even if application corruption were a hundred times more likely, the most likely diagnosis remains unchanged with probability~$1$.
\item \prop{prob\_delta}: \emph{Sensitivity analysis.}\\
  \bprint{P(Problem1=No\_Output)[PrtCbl=Loose↦0.5]−P(Problem1=No\_Output)}\\
  The difference between the failure probability after raising the loose-cable incidence to 50\,\% and the nominal marginal; increase of about 17\%.
\item \prop{prob\_offline}: \emph{Verification under parameter changes.}\\
  \bprint{P(Problem1=No\_Output|PTROFFLINE=Offline∨¬(PrtOn=Yes))}\\
  \bprint{[PrtCbl=Loose↦0.5][PrtMem=Less\_than\_2Mb↦0.4]≥0.6}\\
  Even under simultaneous what-if increases in loose-cable and low-memory incidence, an offline or unpowered printer produces no output with probability at least 0.6.
\end{itemize}
